\documentclass{report}
\usepackage{amsmath,amssymb}
\setlength{\parindent}{0mm}
\setlength{\parskip}{1em}
\begin{document}
\begin{center}
\rule{6in}{1pt} \
{\large Atthaphon Ariyarit \\
{\bf Hybrid Surrogate Model based on Multi-fidelity Efficiency Global Optimization Applied to Helicopter Blade Design}}

Department of Aerospace Engineering \\ Graduate School of System Design \\ Tokyo Metropolitan University \\ 6-6 Asahigaoka \\ Hino-shi \\ Tokyo \\ 191-0065 \\ Japan
\\
{\tt ariyarit-atthaphon@ed.tmu.ac.jp}\\
Masahiko Sugiura\\
Yasutada Tanabe\\
Masahiro Kanzaki\end{center}

In helicopter blades design, it is a problem that computational cost
should be expensive with a high-fidelity computational fluid dynamics
(CFD). Efficient global optimization (EGO) which is Kriging model based
exploration can be reduced the number of expensive function. However, the
reduction of computational cost with CFD is not enough to practical use.
Thus, multi-fidelity approach which combines high-fidelity data and
low-fidelity data has possibility to improve efficiency by reduce the
function of a high-fidelity CFD drastically.

In this study, hybrid Kriging model for multi-fidelity approach are
proposed. In the proposed approach, radial basis function (RBF) is used
to estimate a global model $\mu(x)$. Then, a Kriging method is applied to
evaluate the local deviation $\varepsilon(x)$ ($x$ is set of design
variables). $\mu(x)$ is constructed for data based on a low-fidelity
evalution and $\varepsilon(x)$ is obtained based on a high-fidelity CFD.
In general, the ordinary Kriging model expresses the unknown function
$y(x)=\mu(x)+\varepsilon(x)$ where $\varepsilon(x)$ is a constant. The
proposed multi-fidelity approach is investigated with solving test
problems. Reults by proposed method are compared with single-fidelity
evaluation based surrogate model by an ordinary Kriging method and a
multi-fidelity evaluations based surrogate model by co-Kriging method.
According to this investigations, proposed method can be obtained better
solution than existent method such as an ordinary Kriging based method
and a co-Kriging method. In addition, proposed method should have the
highest efficient method compared with existent methods.

Finally, proposed method is employed to design optimization of helicopter
blades in hovering to obtain the maximum blade efficiency that can be
measured by figure of merit (FOM). Helicopter blades shape are design
through changing their twist distribution. A blade element method (BEM)
solver is used to obtain a low-fidelity data and the computational time
per design is about 0.1 seconds. Reynolds-averaged Navier-Stokes
simulation (RANS) is applied as a high fidelity CFD, which is the
computational time per design is about 48 hours. The optimal shape by the
proposed method could be acquired optimal design which is similar
performance to the optimal design by the high-fidelity evaluation based
optimization. Comparing the proposed method, an ordinary Kriging
single-fidelity approach and a co-Kriging based multi-fidelity approach,
the total number of high-fidelity CFD run was fewest. Thus, the proposed
method which is based on hybrid model is the most efficient optimization
when the optimal design shows similar performance to design by a
single-fidelity based approach.


\end{document}
